\Needspace{8\baselineskip}\originalchapter{2006 作业三: 硬球与阿尔法热化}
题源: 28 Sep 2006, due 5 Oct. 题中DT离子平均质量应为2.5mp; OCR的5.2倒序不采用.
\begin{exercise}\label{ex:06-3-1}
原题1. 两个同半径R硬球碰撞, 入射V0沿x、冲量参数b. (a)靶无限质量, 求Vf; (b)两质量m1,m2且靶初静, 重求; (c)密度n的m2静球海, 入射m1束, 求 $\dot{\mathbf V}_0=-\nu_{sd}\mathbf V_0$ 的初始减速频率.
\end{exercise}
\begin{solution}
接触距离d=2R, b<d. 接触法向n̂取使入射点n̂x为正的等效法线, $\hat n_x=\sqrt{1-b^2/d^2}$; 反弹 $\mathbf V_f=V_0[\mathbf e_x-2\hat n_x\hat{\mathbf n}]$. 有限质量把2改成 $2m_2/(m_1+m_2)$, 切向分量不变. b\ensuremath{\ge}d不碰撞. 乘碰撞率 $nV_0\,2\pi b\dd b$, 仅x平均非零, $\int_0^d2\pi b(1-b^2/d^2)\dd b=\pi d^2/2$, 得 $\nu_{sd}=nV_0\pi d^2m_2/(m_1+m_2)$. 这是束有向动量初始损失而非每碰撞全部能量丢失.
\end{solution}
\begin{exercise}\label{ex:06-3-2}
原题2. 质量mb、电荷Zbe快离子束 $\dot v_b=-(\nu_{be}+\nu_{bi})v_b$, 给 $\nu_{bi}=.094n_{20}/U_b^{3/2}$s$^{-1}$ (U以MeV), $\nu_{be}=100n_{20}/T_e^{3/2}$s$^{-1}$ (Te以keV), ln\ensuremath{\Lambda}20. 为3.5MeV \ensuremath{\alpha}在15keV等密度DT, mb=4mp, mi=2.5mp. (a)求相等频率临界速/能量及 $U_c/T_e$; (b)解vb(t), n20=2, 求tcrit与降到Ub\ensuremath{\simeq}Ti所余\ensuremath{\Delta}t; (c)积分 $\dot U_e=2\nu_{be}U_b$, $\dot U_i=2\nu_{bi}U_b$, 算fe=Ue/(Ue+Ui)、fi, 可数值或解析近似.
\end{exercise}
\begin{solution}
一般快/慢靶频率相等得 $v_c^3=(3\sqrt{2\pi}/(4\pi))T_e^{3/2}/(m_rm_e^{1/2})$, 其中束-离子约化质量 $m_r=m_bm_i/(m_b+m_i)$. 因而 $U_c/T_e=\tfrac{m_b}{2}[3\sqrt{2\pi}/(4\pi m_rm_e^{1/2})]^{2/3}$, 给质量约9.6. 用题给数值系数 $U_c=(.094/100)^{2/3}T_{e,\rm keV}$MeV=.14394MeV.

令a=\ensuremath{\nu}be、b=.094n20, 则 $\dot U=-2(a+bU^{-3/2})U$,
\begin{equation}
 U(t)^{3/2}=(U_0^{3/2}+b/a)\e^{-3at}-b/a,\quad
 t(U)=\frac1{3a}\ln\frac{U_0^{3/2}+b/a}{U^{3/2}+b/a}.
\end{equation}
$v_b(t)=\sqrt{2U(t)/m_b}$且单位换成SI. 实算tcrit=.39716s, tfinal=.46107s, \ensuremath{\Delta}t=.06391s. 到热能附近快离子假设失效, 不将负U的形式继续延拓.

转移能量不应仅用先电子后离子的阶跃分类. 实际 $U_e=\int_{U_f}^{U_0}U^{3/2}[U^{3/2}+U_c^{3/2}]^{-1}\dd U$, $U_i=U_0-U_f-U_e$. Simpson积分得 $f_e=.921079$, $f_i=.078921$; 相对于\textbf{初能}分别.917132、.078583, 另有.004286残留. 分母按题为Ue+Ui, 区分残留才不会混报92.1\%和91.7\%.
\end{solution}
